
\documentclass[10pt,landscape,fleqn]{article}
%
\usepackage{pgfplots}
\usepackage[T1]{fontenc}
\usepackage[utf8]{inputenc}
\usepackage[ngerman]{babel}
\usepackage{lmodern}
\usepackage{pgfplots}
\pgfplotsset{compat=1.12}
\usepgfplotslibrary{fillbetween}
\usetikzlibrary{patterns}
%
%\usepackage{pstricks,pst-plot,pst-text,pst-solides3d}
\usepackage{epsfig}
\usepackage{graphicx}
\usepackage[colorlinks=true,
            pdftitle={Mathe II/B},
            pdfauthor={Thomas Carraro},
            pdfpagemode={UseOutlines},pdfstartview={Fit},
            linkcolor=black,urlcolor=blue,
            pdfsubject={Mathematik IV}]{hyperref}
\usepackage{cancel}
%
\usepackage{ifthen}
\newboolean{mitLoes}
%\setboolean{mitLoes}{false}
\setboolean{mitLoes}{true}

\newcounter{sol}
\setcounter{sol}{0}
\ifthenelse{\boolean{mitLoes}}{
\setcounter{sol}{1}
}

\usepackage{amssymb,amsmath,amsthm,am1-aufgaben,bm}
\usepackage{array, hhline}
\usepackage{multirow, multicol}
\usepackage{paralist}
\usepackage{mathtools}

% \usepackage{marvosym,pifont}
% \usepackage{listings}
% \usepackage{algorithmic,algorithm}
% \lstset{language=Matlab,numbers=left}

\newcounter{Aufg}
\setcounter{Aufg}{0}

\newcommand{\Aufgabe}[3][]{
\stepcounter{Aufg}
\subsubsection*{Aufgabe
\arabic{Aufg}}
{#3}
}

\newcommand{\Loesung}[1]{
	\ifthenelse{\boolean{mitLoes}}
	{\subsubsection*{L\"osung:}
		#1}
	{}
	\ifthenelse{\boolean{mitLoes}}{
	\leavevmode
	\newline
	\rule{360pt}{1pt}
\cleardoublepage 
}

}


\textwidth = 270mm
\textheight = 180mm
\leftmargin 5mm
\oddsidemargin = -20mm
\evensidemargin = -20mm
\topmargin = -25mm
\parindent 0cm
\columnsep 2cm

\newcommand{\ds}{\displaystyle}
\newcommand{\rme}{\mathrm{e}}
\newcommand{\rmi}{\mathrm{i}}

\renewcommand{\vec}{\boldsymbol}
\newcommand{\mat}{\boldsymbol}

\newcommand{\ud}{\;\mathrm{d}}
\newcommand{\p}{\partial}
\newcommand{\dO}{\d \vec o}
\newcommand{\dA}{\d \vec A}
\newcommand{\dX}{\d \vec x}

\renewcommand{\div}{\operatorname{div}}
\newcommand{\rot}{\operatorname{rot}}
\newcommand{\grad}{\operatorname{grad}}

\newcommand{\skalar}[1]{\left\langle{#1}\right\rangle}

\renewcommand{\d}{\,{\rm d}}
\newcommand{\x}{\ensuremath{{\vec x}}}
\newcommand{\ve}[3]{\begin{pmatrix} #1\\ #2\\ #3\end{pmatrix}}
\newcommand{\e}{\ensuremath{{\bf e}}}

\begin{document}



\begin{twocolumn}



\href{http://mb-mathe.hsu-hh.de/bause}
{\large\bf Helmut-Schmidt-Universit\"at}\\
\href{http://mb-mathe.hsu-hh.de/bause}
{\large\bf Universit\"at der Bundeswehr Hamburg}\\
\href{http://mb-mathe.hsu-hh.de/bause}
{\large\bf Fakult\"at f\"ur Maschinenbau und\,Bauingenieurwesen}\\[1ex]
\href{http://mb-mathe.hsu-hh.de/bause}
{\large\bf Prof. Dr.\ Thomas Carraro}\\
\\[-28mm]
\phantom{.}
\hfill\includegraphics[height=20mm]{HSU_RGB}

\vspace*{8mm}
{\Large\bf Mathematik II/B (WI/ET) \hfill Zusatzblatt}\\[1ex]
{\Large WT 2025 \hfill  Taylor-Entwicklung}


\vspace*{1mm}
\rule{360pt}{2pt}

\vspace*{0ex}
% \paragraph*{Ex1}
%
\Aufgabe{}{
Gegeben sei die Funktion $$
f(x) = \frac{1}{(x+1)^2},
$$
\begin{abc}
\item Bestimmen Sie das Taylor-Polynom 3. Ordnung $T_3(x)$ von $f (x)$ um den
	  Punkt $x_0 = 0$.
\item Geben Sie das Restglied $R_3(x;0)$ an. 
\item Schätzen Sie das Restglied für $|x| < \frac{1}{2}$ ab.
\end{abc}
\textbf{Hinweis:} Es ist nicht notwendig das Restglied optimal abzuschätzen.
}


\Loesung{
\begin{abc}
\item Die Taylor-Entwicklung von \( f(x) \) um \( x_0 = 0 \) ist:
\[
f(x) = f(0) + f'(0)x + \frac{f''(0)}{2!}x^2 + \frac{f'''(0)}{3!}x^3 \cdots
\]

\begin{align*}
f(0)  &= 1 \\
f'(x) &= -\frac{2}{(x+1)^3} \\
f'(0) &= -2 \\
f''(x) &= \frac{6}{(x+1)^4} \\
f''(0) &= 6 \\
f'''(x) &= -\frac{24}{(x+1)^5} \\
f'''(0) &= -24
\end{align*}


Damit ergibt sich das Taylor-Polynom zu
\[
T_3(x) = f(0) + f'(0)x + \frac{f''(0)}{2!}x^2 +  \frac{f'''(0)}{3!}x^3 \\
= 1-2x+3x^2-4x^3.
\]

\item Das Restglied wird wie folgt bestimmt:
\[
R_3(x; 0) = \frac{f^{(4)}(\xi)}{4!}x^4,
\]
wobei $-\frac{1}{2}\leq \xi \leq \frac{1}{2}$ gilt.

\[
f^{(4)}(x) = \frac{120}{(x+1)^6}.
\]

Das Restglied wird damit 
\[
R_3(x; 0) = \frac{f^{(4)}(\xi)}{4!}x^4 = \frac{120}{24(\xi+1)^6}x^4.
\]

\item For \( |x| < \frac{1}{2} \):
\[
|R_3(x; 0)| = \left|\frac{120}{24(\xi+1)^6}x^4\right| = \frac{5}{(\xi+1)^6}x^4
\leq \frac{5}{(\frac{1}{2})^6} \left( \frac{1}{2}\right) ^4 = \frac{5 \cdot 64}{16} = 20.
\]

D.h.
\[
|R_3(x; 0)| \leq 20.
\]

	  
\end{abc}

}

\Aufgabe{}{
	Gegeben sei die Funktion 
	$$
	f(x) = e^{-2x}.
	$$
	\begin{abc}
		\item Bestimmen Sie das Taylor-Polynom 3. Ordnung \( T_3(x) \) von \( f(x) \) um den Punkt \( x_0 = 0 \).
		\item Geben Sie das Restglied \( R_3(x;0) \) an. 
		\item Schätzen Sie das Restglied für \( |x| < \frac{1}{2} \) ab.
	\end{abc}
}

\Loesung{
	\begin{abc}
		\item Die Taylor-Entwicklung von \( f(x) \) um \( x_0 = 0 \) ist:
		\[
		T_3(x) = f(0) + f'(0)x + \frac{f''(0)}{2!}x^2 + \frac{f'''(0)}{3!}x^3 + \dots
		\]
		
		Berechnung der Ableitungen:
		
		\begin{align*}
			f(0) &= 1 \\
			f'(x) &= -2e^{-2x} \quad \Rightarrow \quad f'(0) = -2 \\
			f''(x) &= 4e^{-2x} \quad \Rightarrow \quad f''(0) = 4 \\
			f'''(x) &= -8e^{-2x} \quad \Rightarrow \quad f'''(0) = -8
		\end{align*}
		
		Damit ergibt sich das Taylor-Polynom dritter Ordnung:
		
		\[
		T_3(x) = f(0) + f'(0)x + \frac{f''(0)}{2!}x^2 + \frac{f'''(0)}{3!}x^3
		\]
		
		\[
		T_3(x) = 1 - 2x + 2x^2 - \frac{4}{3}x^3.
		\]
		
		\item Das Restglied wird durch die Lagrange-Formel gegeben:
		
		\[
		R_3(x; 0) = \frac{f^{(4)}(\xi)}{4!}x^4,
		\]
		
		
		
		Die vierte Ableitung von \( f(x) \) ist:
		
		\[
		f^{(4)}(x) = 16e^{-2x}.
		\]
		
		Damit folgt für das Restglied:
		
		\[
		R_3(x; 0) = \frac{16e^{-2\xi}}{24}x^4 = \frac{2}{3}e^{-2\xi} x^4.
		\]
		
		\item Für \( |x| < \frac{1}{2} \) setzen wir die obere Schranke für \( e^{-2\xi} \):
		
		Da \( e^{-2\xi} \) auf \( [-\frac{1}{2}, \frac{1}{2}] \) maximal ist für \( \xi = -\frac{1}{2} \), gilt:
		
		\[
		e^{-2\xi} \leq e.
		\]
		
		Damit ergibt sich die Abschätzung:
		
		\[
		|R_3(x; 0)| \leq \frac{2}{3} e \cdot \left(\frac{1}{2}\right)^4.
		\]
		
		\[
		= \frac{2}{3} e \cdot \frac{1}{16} = \frac{2e}{48} = \frac{e}{24}.
		\]
		
		Also:
		
		\[
		|R_3(x; 0)| \leq \frac{e}{24}.
		\]
		
	\end{abc}
}

\Aufgabe{}
{
Gegeben sei die Funktion
\[\
f(x) = \frac{e^x}{(x+1)^2}.
\]

\begin{enumerate}
    \item Bestimmen Sie das Taylor-Polynom 2. Ordnung \(T_2(x)\) von \(f(x)\) um den Punkt \(x_0 = 0\).
    \item Geben Sie das Restglied \(R_2(x;0)\) an.
    \item Schätzen Sie das Restglied für \( |x| < \frac{1}{3} \) ab.
\end{enumerate}

\textbf{Hinweis:} Es ist nicht notwendig, das Restglied optimal abzuschätzen.
}

\Loesung{
\begin{enumerate}
    \item Die Taylor-Entwicklung von \( f(x) \) um \( x_0 = 0 \) lautet:
    \[
    T_2(x) = f(0) + f'(0)x + \frac{f''(0)}{2!}x^2 + \dots
    \]
    
    Berechnung der Ableitungen:
    
    \begin{align*}
    f(0) &= \frac{e^0}{(0+1)^2} = 1 \\
    f'(x) &= \frac{(x-1)e^x}{(x+1)^3} \quad \Rightarrow \quad f'(0) = -1\\
    f''(x) &= \frac{(x^2 - 2x + 3)e^x}{(x+1)^4} \quad \Rightarrow \quad f''(0) = 3
    \end{align*}
    
    Damit ergibt sich das Taylor-Polynom 2. Ordnung:
    
    \[
    T_2(x) = f(0) + f'(0)x + \frac{f''(0)}{2!}x^2
    = 1 - x + \frac{3}{2}x^2.
    \]
    
    \item Das Restglied ist gegeben durch:
    
    \[
    R_2(x; 0) = \frac{f^{(3)}(\xi)}{3!}x^3,
    \]
    
    wobei \( -\frac{1}{3} \leq \xi \leq \frac{1}{3} \).
    
    Dritte Ableitung:
    
    \[
    f'''(x) = \frac{(x^3 - 3x^2 + 9x - 11)e^x}{(x+1)^5}.
    \]
    
    \item Abschätzung des Restglieds für \( |x| < \frac{1}{3} \):
    
    \[
    |R_2(x; 0)| = \left|\frac{f^{(3)}(\xi)}{3!}x^3\right| = \left|\frac{(\xi^3 - 3\xi^2 + 9\xi - 11)e^\xi}{6(\xi+1)^5} x^3\right|.
    \]
    
    Da \( |x| < \frac{1}{3} \), maximieren wir den Bruch für \( -\frac{1}{3} \leq \xi \leq \frac{1}{3} \). Eine grobe Abschätzung liefert:
    \[
    \left| f'''(\xi) \right| = \left| \frac{(\xi^3 - 3\xi^2 + 9\xi - 11)e^\xi}{(\xi+1)^5} \right|
    \]
    
    \[
    \left| f'''(\xi) \right| \leq \frac{\left| \xi^3 - 3\xi^2 + 9\xi - 11 \right| e^\xi}{|(\xi+1)^5|}
    \]
    
    \[
    \left| \xi^3 - 3\xi^2 + 9\xi - 11 \right| \leq \left| \xi^3 \right| + \left| -3\xi^2 \right| + \left| 9\xi \right| + \left| -11 \right|
    \]
    
    Wir w\"ahlen \( \xi = \frac{1}{3} \):
    
    \[
    \frac{1}{27} + \frac{1}{3} + 3 + 11 = \frac{388}{27}
    \]
    
    \[
    (\xi+1)^5 \bigg|_{\xi = -\frac{1}{3}} =  \left(\frac{-1}{3}+1\right)^5=\left(\frac{2}{3}\right)^5 = \frac{32}{243} \approx 0.1317.
    \]
    
    \[
    \max_{\xi} |f'''(\xi)| = \frac{\frac{388}{27} \cdot e^\xi}{(\xi+1)^5} = \frac{\frac{388}{27} \times 1.3956}{1024/243} = 4.76.
    \]
    
    \[
    \left| R_2(x; 0) \right| \leq \frac{4.76}{6} x^3 = 0.79 x^3.
    \]
    
    
    Daraus ergibt sich die Schranke für das Restglied:
    
    \[
    |R_2(x; 0)| \leq 0.79 \left(\frac{1}{3}\right)^3.
    \]
    
\end{enumerate}
}

%\Aufgabe{}{
Gegeben sei die Funktion $$
f(x) = \frac{1}{\sqrt{x^2+1}},
$$
\begin{abc}
\item Bestimmen Sie das Taylor-Polynom 2.Ordnung $T_2(x)$ von $f (x)$ um den
	  Punkt $x_0 = 1$.
\item Geben Sie das Restglied $R_2(x;1)$ an.
\item Schätzen Sie das Restglied für $\frac{1}{2} < x < \frac{3}{2}$ ab.
\end{abc}
\textbf{Hinweis:} Es ist nicht notwendig das Restglied optimal abzuschätzen.
}


\Loesung{
\begin{abc}
\item Die Taylor-Entwicklung von \( f(x) \) um \( x_0 = 1\) ist:
\[
f(x) = f(1) + f'(1)(x-1) + \frac{f''(1)}{2!}(x-1)^2 \cdots
\]

\begin{align*}
f(1)  &= \frac{1}{\sqrt{2}} \\
f'(x) &= -\frac{x}{\sqrt{(x^2+1)^3}}  \quad \Rightarrow \quad f'(1) = -\frac{1}{2\sqrt{2}}\\
f''(x) &= \frac{2x^2-1}{\sqrt{(x^2+1)^5}}  \quad \Rightarrow \quad f''(1) = \frac{1}{4\sqrt{2}}
\end{align*}


Damit ergibt sich das Taylor-Polynom zu
\[
T_2(x) = f(1) + f'(1)(x-1) + \frac{f''(1)}{2!}(x-1)^2 \\
= \frac{1}{\sqrt{2}}-\frac{1}{2\sqrt{2}}(x-1)+\frac{1}{8\sqrt{2}}(x-1)^2.
\]

\item Das Restglied wird wie folgt bestimmt:
\[
R_2(x; 1) = \frac{f'''(\xi)}{3!}(x-1)^3,
\]
wobei $\frac{1}{2}\leq \xi \leq \frac{3}{2}$ gilt.

\[
f'''(x) = -\frac{6x^3-9x}{\sqrt{(x^2+1)^7}}.
\]

Das Restglied wird damit 
\[
R_2(x; 1) = \frac{f'''(\xi)}{3!}(x-1)^3
=-\frac{2\xi^3-3\xi}{2\sqrt{(\xi^2+1)^7}}(x-1)^3.
\]


\item F\"ur \( \frac{1}{2} <x  < \frac{3}{2} \) ergibt sich \( \frac{1}{2} \leq \xi \leq \frac{3}{2} \) :
\begin{align*}
|R_2(x; 1)|& = \left|-\frac{2\xi^3-3\xi}{2\sqrt{(\xi^2+1)^7}}(x-1)^3\right|\\
&=\left|-\frac{\xi(2\xi^2-3)}{2\sqrt{(\xi^2+1)^7}}(x-1)^3\right|\\
&\leq \left|-\frac{\frac{3}{2}(2\frac{3}{2}^2-3)}{2\sqrt{(\frac{1}{2}^2+1)^7}}(\frac{3}{2}-1)^3\right|\\
&= \frac{144}{\sqrt{5^7}}\frac{1}{8}\\
&= \frac{18}{\sqrt{5^7}}\\
&\leq \frac{18}{\sqrt{4^7}}\\
&= \frac{18}{2^7}\\
&= \frac{9}{2^6}  
\end{align*}

D.h.
\[
|R_2(x; 1)| \leq \frac{9}{2^6} .
\]

	  
\end{abc}

}

\Aufgabe{}{
Gegeben sei die Funktion $$
f(x) = \cos(x^2),
$$
\begin{abc}
\item Bestimmen Sie das Taylor-Polynom 3. Ordnung $T_3(x)$ von $f (x)$ um den
	  Punkt $x_0 = 0$.
\item Geben Sie das Restglied $R_3(x;0)$ an. 
\item Schätzen Sie das Restglied für $|x| < \frac{1}{2}$ ab.
\end{abc}
\textbf{Hinweis:} Es ist nicht notwendig das Restglied optimal abzuschätzen.
}


\Loesung{

	
	\begin{abc}
		\item
	
	Das Taylor-Polynom \( n \)-ter Ordnung einer Funktion \( f(x) \) um den Punkt \( x_0 \) ist gegeben durch:
	
	\[
	T_n(x) = \sum_{k=0}^n \frac{f^{(k)}(x_0)}{k!} (x - x_0)^k
	\]
	
	Für \( n = 3 \) und \( x_0 = 0 \) lautet das Taylor-Polynom:
	
	\[
	T_3(x) = f(0) + f'(0) \cdot x + \frac{f''(0)}{2!} \cdot x^2 + \frac{f'''(0)}{3!} \cdot x^3
	\]
	
	\subsubsection*{Schritt 1: Berechnung der Ableitungen von \( f(x) = \cos(x^2) \)}
	
	\begin{align*}
	1. \quad & f(x) = \cos(x^2) \\
	2. \quad & f'(x) = -\sin(x^2) \cdot 2x \\
	3. \quad & f''(x) = -\cos(x^2) \cdot (2x)^2 - \sin(x^2) \cdot 2 = -4x^2 \cos(x^2) - 2 \sin(x^2) \\
	4. \quad & f'''(x) = 8x^3 \sin(x^2) - 12x \cos(x^2)
	\end{align*}
	
	\subsubsection*{Schritt 2: Auswertung der Ableitungen an der Stelle \( x_0 = 0 \)}
	
	\begin{align*}
	1. \quad & f(0) = \cos(0) = 1 \\
	2. \quad & f'(0) = -\sin(0) \cdot 0 = 0 \\
	3. \quad & f''(0) = -4 \cdot 0^2 \cos(0) - 2 \sin(0) = 0 \\
	4. \quad & f'''(0) = 8 \cdot 0^3 \sin(0) - 12 \cdot 0 \cos(0) = 0
	\end{align*}
	
	\subsubsection*{Schritt 3: Einsetzen in das Taylor-Polynom}
	
	\[
	T_3(x) = 1 + 0 \cdot x + \frac{0}{2!} \cdot x^2 + \frac{0}{3!} \cdot x^3 = 1
	\]
	
	Das Taylor-Polynom 3. Ordnung von \( f(x) = \cos(x^2) \) um \( x_0 = 0 \) lautet also:
	
	\[
	T_3(x) = 1
	\]
	
		
	\item Abschätzung des Restglieds für \( |x| < \frac{1}{2} \)
		
		Wir schätzen das Restglied \( R_3(x;0) \) für \( |x| < \frac{1}{2} \) ab. Dazu betrachten wir den Betrag des Restglieds:
		
		\[
		|R_3(x;0)| = \left| \frac{f^{(4)}(\xi)}{24} x^4 \right|
		\]
		
		Da \( |x| < \frac{1}{2} \), gilt \( x^4 < \left( \frac{1}{2} \right)^4 = \frac{1}{16} \). Um \( |f^{(4)}(\xi)| \) abzuschätzen, betrachten wir den maximalen Wert der 4. Ableitung im Intervall \( |\xi| < \frac{1}{2} \).
		
		\subsubsection*{Schritt 1: Abschätzung von \( |f^{(4)}(\xi)| \)}
		
		Für \( |\xi| < \frac{1}{2} \):
		
		\[
		|f^{(4)}(\xi)| = \left| 16\xi^4 \cos(\xi^2) - 48\xi^2 \sin(\xi^2) - 12 \cos(\xi^2) \right|
		\]
		
		Da \( |\cos(\xi^2)| \leq 1 \) und \( |\sin(\xi^2)| \leq 1 \), gilt:
		
		\[
		|f^{(4)}(\xi)| \leq 16 \cdot \left( \frac{1}{2} \right)^4 \cdot 1 + 48 \cdot \left( \frac{1}{2} \right)^2 \cdot 1 + 12 \cdot 1 = 16 \cdot \frac{1}{16} + 48 \cdot \frac{1}{4} + 12 = 1 + 12 + 12 = 25
		\]
		
		\subsubsection*{Schritt 2: Abschätzung des Restglieds}
		
		\[
		|R_3(x;0)| \leq \frac{25}{24} \cdot \frac{1}{16} = \frac{25}{384} \approx 0,0651
		\]
		
		Für \( |x| < \frac{1}{2} \) gilt also:
		
		\[
		|R_3(x;0)| \leq \frac{25}{384}
		\]
		
		
	\end{abc}
	


}

% 



%%%%%%%%%%%%%%%%%%%%%%%%%%%%%%%%%%%%%%%%%%%%%%%%%%%%%%%%%%%%%%%%%%%%%%%%%%%%%%%%


%%%%%%%%%%%%%%%%%%%%%%%%%%%%%%%%%%%%%%%%%%%%%%%%%%%%%%%%%%%%%%%%%%%%%%%%%%%%%%%%
\end{twocolumn}

\end{document}

%%%%%%%%%%%%%%%%%%%%%%%%%%%%%%%%%%%%%%%%%%%%%%%%%%%%%%%%%%%%%%%%%%%%%%%%%%%%%%%%
